% Conference poster, 140 x 90 cm landscape (NeurIPS main-conference size), makeposter house style.
% Every coordinate is in centimetres on the printed sheet, with the origin at the bottom left.
% Three columns (39, 56 and 39 cm) over a full-width row of plots; reading order is left column,
% middle, right column, then the plots.
% Build:  python3 ~/.claude/skills/makeposter/scripts/poster.py build     (or: pdflatex poster.tex)
% Keep box coordinates as plain numbers: poster.py reads them to crop, measure and proofread.
\documentclass[tikz,border=0pt]{standalone}
\usepackage{posterstyle}
\newcommand{\PW}{140}\newcommand{\PH}{90}   % page size in cm

\begin{document}
\begin{tikzpicture}[x=1cm,y=1cm]
\useasboundingbox (0,0) rectangle (\PW,\PH);
\fill[white] (0,0) rectangle (\PW,\PH);

% =================================== HEADER ===================================================
\posterband{25}{78.4}{117}{88.6}{5.5}
  {Short Name of the Method}
  {The Rest of the Paper Title, Set as the Subtitle}
  {First Author, Second Author, Third Author, and Last Author}
% left: venue logo over institution logo (vector files; a placeholder is drawn while a file is missing)
\node[inner sep=0pt] at (13,85.5) {\posterimg{20cm}{5.6cm}{logos/venue_logo.pdf}};
\node[inner sep=0pt] at (13,80.2) {\posterimg{20cm}{2.4cm}{logos/institution_logo.pdf}};
% right: QR codes, top edge level with the band
\posterqr{123.3}{88.6}{7.6cm}{https://arxiv.org/abs/0000.00000}{Paper}
\posterqr{133.4}{88.6}{7.6cm}{https://github.com/user/repository}{Code}

% =================================== LEFT COLUMN ==============================================
\pbox{1.5}{58.85}{40.5}{77.0}{The Problem, in a Noun Phrase}{%
Two sentences that state the problem for a visitor from another area.
Say what current methods need and where that breaks.\par
\vspace{0.6cm}
{\centering\posterimg{30cm}{5.2cm}{figures/problem_sketch.pdf}\par}
\vspace{0.6cm}
One sentence that lists the consequences, in the words the paper uses.
}

\pbox{1.5}{30.5}{40.5}{56.95}{Formal Setting}{%
\hl{Definition (Name).} One sentence that says what is measured, then
\[
\phi(r)=1-|\tau(r,a)|,\qquad \tau(r,a)=\frac{N_{+}(r)-N_{-}(r)}{\binom{n}{2}},
\]
where every symbol is explained in words, with the \hl{new terms} in bold once.\par
\vspace{0.3cm}
\hl{Definition (Second name).} The second quantity that the problem needs,
\[
d(r_0,r)=\sum_{(i,j)\ \text{inverted}} w_{ij},
\]
with its terms named here, because later boxes and plots reuse exactly these names.\par
\vspace{0.3cm}
\hl{Problem.} Given a threshold $\psi\in[0,1]$, find
\[
\min_{r}\; d(r_0,r)\quad\text{subject to}\quad \phi(r)\ge\psi .
\]
}

% =================================== MIDDLE COLUMN ============================================
\pbox{42}{30.5}{98}{77.0}{Our Solution: Method Name (ACRONYM)}{%
Two sentences with the one idea behind the method, the \hl{key term} in bold.
Say what a single step changes, in the quantities defined on the left.\par
\vspace{1.4cm}
\noindent\begin{posteralg}[0.57\linewidth]{Method Name (ACRONYM)}
1: & $r\gets r_0$\\
2: & \textbf{while} $\phi(r)<\psi$ and a step can raise $\phi(r)$ \textbf{do}\\
3: & \quad apply the cheapest such step\\
4: & \textbf{return} $r$
\end{posteralg}\hfill
\begin{minipage}[c]{0.39\linewidth}
\raggedright
Two sentences next to the algorithm: the data structure behind line 3 and what one step costs.
\end{minipage}\par
\vspace{1.9cm}
{\centering\posterimg{\linewidth}{22.3cm}{figures/fig_example.pdf}\par}
\vspace{1.4cm}
{\fCaption Two lines that say how to read the example (what the colours mean) and give its outcome in
numbers that a visitor can check against the figure.\par}
}

% =================================== RIGHT COLUMN =============================================
\pbox{99.5}{52.9}{138.5}{77.0}{The Main Result in Four Words}{%
\hl{Definition (Name).} The one extra quantity the theorem needs, with its formula:
\[
\Delta(r)=|N_{+}(r)-N_{-}(r)| .
\]\par
\vspace{0.6cm}
\hl{Theorem.} The statement as a single sentence that ends in
\[
K^{\ast}=\left\lceil\frac{\Delta_0-\Delta^{\ast}}{2}\right\rceil\ \ \text{steps, at distance}\ \
d(r_0,r_{\mathrm{out}})=\frac{K^{\ast}}{\binom{n}{2}}.
\]\par
\vspace{0.6cm}
\hl{Proof idea.} One sentence for the lower bound that every solution obeys.
One sentence for why the method attains it.
}

\pbox{99.5}{30.5}{138.5}{51.0}{Guarantees and Runtime}{%
\noindent\begin{tabular*}{\linewidth}{@{}l@{\extracolsep{\fill}}ll@{}}
\toprule
\textbf{Setting} & \textbf{Guarantee} & \textbf{Time}\\
\midrule
first & optimal & $O(n\log n+K^{\ast})$\\[0.25cm]
second & $c$-approximation & $O((n+K^{\ast})\log n)$\\
\bottomrule
\end{tabular*}\par
\vspace{0.3cm}
{\fCaption One line under the table for what both rows share and for symbols not yet defined.\par}
\vspace{0.6cm}
Two sentences for results that do not fit a table cell: a special case with its exact setting,
and any conjecture in the paper's own words.
}

% =================================== BOTTOM ROW ===============================================
% Titles hang from y=27.2, the plots stand on y=6.9 and the captions start at y=5.8.
% Give the plots one common height, so that their x-axes sit on one line.
\pframe{1.5}{1.4}{138.5}{29.0}
\ppanelsetup{27.2}{6.9}{5.8}
\ppanel{34.5}{59}{The main empirical claim, as a full sentence}
  {\posterimg{60cm}{17.9cm}{figures/fig_main_poster.pdf}}
  {The first sentence says what each mark is. The second says how to read the plot,
   for example which direction is better.}
\ppanel{82}{29}{A second claim that a plot supports}
  {\posterimg{29cm}{17.15cm}{figures/fig_second_poster.pdf}}
  {The first sentence gives the setting of the plot. The second gives the reason behind what it shows.}
\ppanel{117}{36}{The method scales to large inputs}
  {\posterimg{36.9cm}{17.15cm}{figures/fig_runtime_poster.pdf}}
  {One sentence with the headline number. One sentence on the hardest case.}
\end{tikzpicture}
\end{document}
